\uuid{ZOZR}
\exo7id{4689}
\titre{exo7 4689}
\auteur{quercia}
\organisation{exo7}
\datecreate{2010-03-16}
\isIndication{false}
\isCorrection{true}
\chapitre{Suite}
\sousChapitre{Convergence}
\cpge{mpsi-mp2i.nombres-reels-suites.b, pcsi.nombres-reels-suites.a, ptsi.nombres-reels-suites.a, mpsi-mp2i.nombres-reels-suites.c, pcsi.nombres-reels-suites.b, ptsi.nombres-reels-suites.b, mpsi-mp2i.nombres-reels-suites.d, pcsi.nombres-reels-suites.c, ptsi.nombres-reels-suites.c, bcpst1.suites-reelles}
\module{Analyse}
\niveau{L1}
\difficulte{}

\contenu{
\texte{
Pour $n \in \N^*$, on note $S(n)$ la somme des chiffres de l'{\'e}criture d{\'e}cimale de $n$.
}
\begin{enumerate}
    \item \question{Encadrer $S(n+1)$ en fonction de $S(n)$.
    En d{\'e}duire que la suite $\left(\frac {S(n+1)}{S(n)}\right)$ est born{\'e}e.}
    \item \question{Chercher $\inf\left\{\frac {S(n+1)}{S(n)} \text{ tq } n \in \N^*\right\}$, et
             $\sup\left\{\frac {S(n+1)}{S(n)} \text{ tq } n \in \N^*\right\}$.}
    \item \question{La suite $\left(\frac {S(n+1)}{S(n)}\right)$ est-elle convergente ?}
\reponse{
$1 \le S(n+1) \le S(n) + 1  \Rightarrow  0 \le \frac {S(n+1)}{S(n)} \le 2$.
$\inf = 0$ $(99\dots99)$, $\sup = 2$ $(100\dots00)$.
}
\end{enumerate}
}
